\documentclass{beamer}

\usepackage{Archivo}
\renewcommand{\sfdefault}{\rmdefault}
\usepackage{graphicx}
\usepackage{tikz}
\usetikzlibrary{calc}
\usepackage[most]{tcolorbox}

%% ---- Superteam Poland brand colors ----
\definecolor{stpldark}{HTML}{0B0E0F}   % near-black background
\definecolor{stplcard}{HTML}{161A1B}   % card / block surface
\definecolor{stplborder}{HTML}{1E2526}  % subtle track color for progress bar
\definecolor{stplred}{HTML}{DC143C}    % brand crimson red
\definecolor{stplwhite}{HTML}{FFFFFF}  % white
\definecolor{stplmuted}{HTML}{8A9BA8}  % muted / secondary text

\usetheme{default}
\setbeamersize{text margin left=0.85cm, text margin right=0.85cm}

%% ---- Background & base text ----
\setbeamercolor{background canvas}{bg=stpldark}
\setbeamercolor{normal text}{fg=stplwhite, bg=stpldark}
\usebeamercolor{normal text}

%% ---- Background image helpers ----
\newcommand{\stpldefaultbg}{%
  \setbeamertemplate{background}{%
    \begin{tikzpicture}[remember picture, overlay]
      \node[at=(current page.center), inner sep=0pt] {%
        
\includegraphics[width=\paperwidth,height=\paperheight]{images/bg-default.png}%
      };
      \fill[stplborder]
        (current page.north west) rectangle
        ([yshift=-2.5pt]current page.north east);
      \pgfmathsetmacro{\stplprog}{\insertframenumber/\inserttotalframenumber}
      \fill[stplred]
        (current page.north west) rectangle
        ($(current page.north west)+(\stplprog\paperwidth,-2.5pt)$);
    \end{tikzpicture}%
  }%
}
\newcommand{\stplherobg}{%
  \setbeamertemplate{background}{%
    
\includegraphics[width=\paperwidth,height=\paperheight]{images/bg-hero.png}%
  }%
}
\stpldefaultbg

%% ---- Structure ----
\setbeamercolor{structure}{fg=stplred}

%% ---- Title slide metadata ----
\setbeamercolor{title}{fg=stplwhite}
\setbeamercolor{subtitle}{fg=stplred}
\setbeamercolor{author}{fg=stplwhite}
\setbeamercolor{institute}{fg=stplmuted}
\setbeamercolor{date}{fg=stplmuted}

%% ---- Frame title: flat, no box background ----
\setbeamercolor{frametitle}{fg=stplwhite, bg=}
\setbeamerfont{frametitle}{size=\large, series=\bfseries}
\setbeamertemplate{frametitle}{%
  \vspace{0.35cm}%
  {\centering\insertframetitle\par}%
  \nointerlineskip\vspace{0.18cm}%
  \hfil{\color{stplred}\rule{2.2cm}{1.5pt}}\hfil\par%
}



%% ---- Incremental reveals ----
\beamerdefaultoverlayspecification{<+->}

%% ---- Bullet items ----
\setbeamercolor{itemize item}{fg=stplred}
\setbeamercolor{itemize subitem}{fg=stplmuted}
\setbeamercolor{itemize subsubitem}{fg=stplmuted}
\setbeamertemplate{itemize item}{\bfseries\textendash}
\setbeamertemplate{itemize subitem}{\small\textendash}
\setbeamercolor{enumerate item}{fg=stplred}
\setbeamercolor{enumerate subitem}{fg=stplmuted}

%% ---- Blocks (flat left-accent via tcolorbox) ----
\tcbset{
  stplbase/.style={
    enhanced, breakable,
    colback=stplcard, coltext=stplwhite,
    coltitle=stplwhite, fonttitle=\bfseries,
    boxrule=0pt, leftrule=3pt,
    arc=0pt, outer arc=0pt,
    top=5pt, bottom=5pt, left=7pt, right=7pt,
    toptitle=4pt, bottomtitle=4pt,
    titlerule=0pt,
  },
}
\renewenvironment{block}[1]{%
  \begin{tcolorbox}[stplbase, colframe=stplred, title={#1}]%
}{%
  \end{tcolorbox}%
}
\renewenvironment{alertblock}[1]{%
  \begin{tcolorbox}[stplbase, colframe=stplred!75!black,
    colback=stplred!10!stplcard, title={#1}]%
}{%
  \end{tcolorbox}%
}
\renewenvironment{exampleblock}[1]{%
  \begin{tcolorbox}[stplbase, colframe=stplmuted, title={#1}]%
}{%
  \end{tcolorbox}%
}

%% ---- Table of contents ----
\setbeamercolor{section in toc}{fg=stplwhite}
\setbeamercolor{subsection in toc}{fg=stplmuted}
\setbeamertemplate{section in toc}{%
  {\color{stplred}\textbf{\textendash}}\hspace{0.5em}\inserttocsection%
}

%% ---- Footline: minimal ----
\setbeamertemplate{footline}{%
  \begin{beamercolorbox}[wd=\paperwidth, ht=0pt, dp=0.55cm]{normal text}%
    \hspace{0.85cm}%
    
\includegraphics[height=1.6ex]{images/stpl-logo.png}%
    \hfill%
    \textcolor{stplmuted}{\tiny\insertframenumber\,/\,\inserttotalframenumber}%
    \hspace{0.85cm}%
  \end{beamercolorbox}%
}
\setbeamertemplate{navigation symbols}{}

%% ---- Title frame helper (uses hero background) ----
\newcommand{\maketitleframe}{%
  \stplherobg
  \begin{frame}[plain]%
    \titlepage%
  \end{frame}%
  \stpldefaultbg
}

%% ---- Section divider slides ----
\AtBeginSection[]{%
  \stplherobg
  \begin{frame}[plain]
    \begin{tikzpicture}[remember picture, overlay]
      \fill[stplred, opacity=0.07]
        ([xshift=-2.5cm]current page.east) circle (5.5cm);
      \fill[stplred]
        (current page.north west) rectangle ([xshift=5pt]current page.south west);
    \end{tikzpicture}%
    \vfill
    \hspace{0.1\paperwidth}%
    \begin{minipage}{0.8\paperwidth}
      \textcolor{stplmuted}{\small Section \thesection}\\[0.35cm]
      {\color{stplwhite}\LARGE\bfseries\insertsectionhead}\\[0.25cm]
      {\color{stplred}\rule{3cm}{1.5pt}}
    \end{minipage}%
    \vfill
  \end{frame}%
  \stpldefaultbg
}

%% ---- Title page ----
\setbeamertemplate{title page}{%
  \begin{tikzpicture}[remember picture, overlay]
    \fill[stplred]
      (current page.north west) rectangle ([xshift=5pt]current page.south west);
    \fill[stplred, opacity=0.07]
      ([xshift=-3cm,yshift=3cm]current page.north east) circle (5.5cm);
    \fill[stplred, opacity=0.04]
      ([yshift=-1.5cm]current page.south east) circle (3.5cm);
  \end{tikzpicture}%
  \vfill
  \hspace{0.08\paperwidth}%
  \begin{minipage}{0.85\paperwidth}
    
\includegraphics[width=3cm]{images/stpl-logo.png}\par
    \vspace{0.8cm}
    {\usebeamerfont{title}\usebeamercolor[fg]{title}\inserttitle}\par
    \vspace{0.2cm}
    {\color{stplred}\rule{3cm}{1.5pt}}\par
    \vspace{0.3cm}
    {\usebeamerfont{subtitle}\usebeamercolor[fg]{subtitle}\insertsubtitle}\par
    \vspace{0.65cm}
    {\usebeamerfont{author}\usebeamercolor[fg]{author}\insertauthor}\par
    \vspace{0.12cm}
    {\small\textcolor{stplmuted}{\insertinstitute\enspace·\enspace\insertdate}}
  \end{minipage}%
  \vfill
}


\usepackage[utf8]{inputenc}
\usepackage[T1]{fontenc}

\title{Pierwszy wallet}
\subtitle{Solana Bootcamp --- Phantom, devnet, Explorer}
\author{}
\institute{Superteam Poland}

% Placeholder for a screenshot to be added later.
% Usage: \screenshotph{opis}  or  \screenshotph[height]{opis}
% Replace with \includegraphics[...]{images/...} once the image is ready.
\newcommand{\screenshotph}[2][3.4cm]{%
  \begin{tikzpicture}
    \node[draw=stplmuted, dashed, line width=0.8pt,
          rounded corners=4pt, fill=stplcard,
          minimum width=\linewidth, minimum height=#1,
          align=center, text=stplmuted, font=\scriptsize\itshape,
          inner sep=4pt]
      {[ zrzut ekranu ]\\#2};
  \end{tikzpicture}}

\begin{document}

\maketitleframe

\begin{frame}{Plan}
  \tableofcontents
\end{frame}

% ================================================================
\section{Klucze}

% TODO: 1. Portfel --- twój login --- hook back to Lekcja 3 (no helper block)
\begin{frame}{Portfel --- twój login}
  \begin{itemize}
    \item Portfel to \textbf{tożsamość} użytkownika w sieci
    \item Ten sam portfel podpisze każdą transakcję do końca kursu
    \item Czemu nie ma hasła?
  \end{itemize}
\end{frame}

% TODO: 2. Para kluczy --- recap (introduced in Lekcja 1)
\begin{frame}{Para kluczy}
  \begin{itemize}
    \item \textbf{Klucz publiczny} pełni rolę adresu --- można go swobodnie udostępniać
    \item \textbf{Klucz prywatny} autoryzuje transakcje --- nigdy nie opuszcza portfela
  \end{itemize}
  \vspace{0.6cm}
  \begin{center}
    \begin{tikzpicture}[
        every node/.style={font=\scriptsize\bfseries},
        box/.style={
          draw=stplred, line width=1pt,
          fill=stplcard, text=stplwhite,
          rounded corners=2pt,
          minimum width=2.3cm, minimum height=0.75cm,
          align=center,
        },
        arr/.style={-latex, draw=stplred, line width=1pt},
        lbl/.style={font=\tiny\itshape, text=stplmuted},
      ]
      % klucz prywatny w portfelu --- wspólny motyw w całej talii
      \node[box] (priv) {Klucz\\prywatny};
      \draw[draw=stplred, dashed, line width=0.8pt, rounded corners=4pt]
        ($(priv.north west)+(-0.22,0.18)$) rectangle ($(priv.south east)+(0.22,-0.18)$);
      \node[font=\tiny\bfseries, text=stplmuted]
        at ($(priv.north)+(0,0.33)$) {Portfel};
      \node[box, right=5.2cm of priv] (pub) {Klucz publiczny\\(adres)};
      \coordinate (wr) at ($(priv.east)+(0.22,0)$);
      \draw[arr] (wr) -- node[lbl, above]{wyprowadza (jednokierunkowo)} (pub.west);
      \node[lbl, below=0.18cm] at ($(wr)!0.5!(pub.west)$)
        {z publicznego nie policzysz prywatnego};
    \end{tikzpicture}
  \end{center}
\end{frame}

% TODO: 3. Podpis --- recap (introduced in Lekcja 1)
\begin{frame}{Podpis}
  \begin{itemize}
    \item Klucz prywatny podpisuje transakcję --- każdy węzeł weryfikuje podpis kluczem publicznym
    \item Zmiana jednego bajtu unieważnia podpis --- treści i autora nie można sfałszować
  \end{itemize}
  \vspace{0.1cm}
  \begin{center}
    \begin{tikzpicture}[
        every node/.style={font=\scriptsize\bfseries},
        box/.style={
          draw=stplred, line width=1pt,
          fill=stplcard, text=stplwhite,
          rounded corners=2pt,
          minimum width=2cm, minimum height=0.65cm,
          align=center,
        },
        arr/.style={-latex, draw=stplred, line width=1pt},
        lbl/.style={font=\tiny\itshape, text=stplmuted},
      ]
      % strona podpisu: portfel z kluczem prywatnym
      \node[box] (priv) {Klucz\\prywatny};
      \draw[draw=stplred, dashed, line width=0.8pt, rounded corners=4pt]
        ($(priv.north west)+(-0.22,0.18)$) rectangle ($(priv.south east)+(0.22,-0.18)$);
      \node[font=\tiny\bfseries, text=stplmuted] at ($(priv.north)+(0,0.33)$) {Portfel};
      \node[box, below=1cm of priv] (pub) {Klucz\\publiczny};
      % para kluczy: publiczny wyprowadzony z prywatnego
      \draw[arr] ($(priv.south)+(0,-0.18)$) -- node[lbl, right=1pt]{wyprowadza} (pub.north);
      \node[box, right=5.2cm of priv] (sig) {Podpis};
      \node[box, right=5.2cm of pub] (node) {Węzeł sieci\\$\checkmark$};
      \coordinate (wr) at ($(priv.east)+(0.22,0)$);
      % podpisanie
      \draw[arr] (wr) -- node[lbl, above]{podpisuje} (sig.west);
      % weryfikacja zbiega dwa wejścia: podpis + klucz publiczny
      \draw[arr] (sig.south) -- node[lbl, right=2pt]{transakcja} (node.north);
      \draw[arr] (pub.east) -- node[lbl, above]{klucz publiczny} (node.west);
    \end{tikzpicture}
  \end{center}
\end{frame}

% ================================================================
\section{Phantom}

% TODO: 4. Instalacja Phantom --- demo (screen-grab moment, no helper block)
\begin{frame}{Instalacja Phantom}
  \begin{columns}[T]
    \column{0.54\textwidth}
      \begin{itemize}
        \item Zainstaluj rozszerzenie \textbf{Phantom} (Chrome, Brave lub Firefox)
        \item Kliknij \textit{,,Create New Wallet''} --- portfel generuje parę kluczy lokalnie
        \item Zapisz \textbf{seed phrase} --- szczegóły na kolejnym slajdzie
        \item Ustaw hasło lokalne --- chroni tylko tę przeglądarkę, nie zastępuje seed
      \end{itemize}
    \column{0.42\textwidth}
      \centering
      \includegraphics[width=\linewidth,height=4.8cm,keepaspectratio]{images/phantom-balance}
  \end{columns}
\end{frame}

% TODO: 5. Seed phrase --- jedna zasada bezpieczeństwa (alertblock: kto ma seed, ma środki)
\begin{frame}{Seed phrase}
  \begin{itemize}
    \item \textbf{Seed phrase} --- 12 lub 24 słowa generowane przy zakładaniu portfela
    \item Z tych słów deterministycznie powstają \textbf{wszystkie} klucze prywatne
    \item Nigdy online --- bez wklejania, bez screenshotów, bez chmury
    \item Zapis offline --- kartka w sejfie lub lokalny manager haseł
  \end{itemize}
  \vspace{0.15cm}
  \begin{center}
    \begin{tikzpicture}[
        every node/.style={font=\scriptsize\bfseries},
        box/.style={
          draw=stplred, line width=1pt,
          fill=stplcard, text=stplwhite,
          rounded corners=2pt,
          minimum width=2cm, minimum height=0.6cm,
          align=center,
        },
        seedbox/.style={
          draw=stplred, line width=1pt,
          fill=stplred, text=stplwhite,
          rounded corners=2pt,
          minimum width=2cm, minimum height=1.2cm,
          align=center,
        },
        arr/.style={-latex, draw=stplred, line width=1pt},
        lbl/.style={font=\tiny\itshape, text=stplmuted},
      ]
      \node[seedbox] (seed) {Seed\\phrase};
      \node[box, right=2.6cm of seed, yshift=0.6cm]  (k1) {Konto 1};
      \node[box, right=2.6cm of seed]                (k2) {Konto 2};
      \node[box, right=2.6cm of seed, yshift=-0.6cm] (k3) {Konto 3};
      \draw[arr] (seed.east) -- (k1.west);
      \draw[arr] (seed.east) -- (k2.west);
      \draw[arr] (seed.east) -- (k3.west);
      \node[lbl, above] at ($(seed.east)!0.5!(k1.west)+(0,0.35cm)$)
        {deterministycznie};
    \end{tikzpicture}
  \end{center}
  \begin{alertblock}{Kto ma seed, ma środki}
    Nie istnieje procedura ,,zapomniałem''. Wycieka raz, na zawsze.
  \end{alertblock}
\end{frame}

% TODO: 6. Devnet vs Mainnet --- A-vs-B opposition (no helper block)
\begin{frame}{Devnet vs Mainnet}
  \begin{columns}
    \column{0.48\textwidth}
      \begin{block}{Devnet}
        \begin{itemize}
          \item Piaskownica dla deweloperów
          \item Tokeny testowe \textbf{bez wartości}
          \item Darmowy airdrop SOL na żądanie
          \item Cały kurs odbywa się tutaj
        \end{itemize}
      \end{block}

    \column{0.48\textwidth}
      \begin{alertblock}{Mainnet}
        \begin{itemize}
          \item Realna sieć produkcyjna
          \item Tokeny mają \textbf{realną wartość}
          \item Każda transakcja kosztuje prawdziwe SOL
          \item Nie używamy w tym kursie
        \end{itemize}
      \end{alertblock}
  \end{columns}
  \vspace{0.3cm}
  \begin{center}
    \textcolor{stplmuted}{Jak nie napisano inaczej --- pracujemy na devnecie.}
  \end{center}
\end{frame}

% TODO: 7. Airdrop --- devnet faucet demo (exampleblock)
\begin{frame}{Airdrop}
  \begin{columns}[T]
    \column{0.54\textwidth}
      \begin{itemize}
        \item Devnet pozwala \textbf{dosypać sobie} testowego SOL z faucetu
        \item Trochę SOL jest potrzebne na opłaty transakcji i \textbf{rent} kont
        \item Źródła --- \texttt{solana airdrop} (CLI), \texttt{faucet.solana.com}, Solana Playground
        \item Typowo 1--2 SOL na żądanie, z dziennym limitem per adres
      \end{itemize}
    \column{0.42\textwidth}
      \includegraphics[width=\linewidth]{images/faucet-devnet-crop}
  \end{columns}
  \vspace{0.05cm}
  \begin{exampleblock}{Demo --- faucet.solana.com}
    Otwórz faucet, wklej adres skopiowany z Phantom, kliknij \textit{,,Devnet''},
    sprawdź zaktualizowane saldo w portfelu.
  \end{exampleblock}
\end{frame}

% TODO: 8. Pierwsza transakcja + Explorer --- demo (exampleblock: podpis/status/opłata)
\begin{frame}{Pierwsza transakcja}
  \begin{itemize}
    \item Wykonaj dowolny transfer SOL między dwoma adresami devnet
    \item Otwórz \texttt{explorer.solana.com} i przełącz sieć na \textbf{Devnet}
    \item Znajdź transakcję po podpisie skopiowanym z Phantom
    \item To dokładnie ten \textbf{write} z Lekcji 1 --- podpisany, opłacony, atomowy
  \end{itemize}
  \vspace{0.15cm}
  \begin{center}
    \begin{tikzpicture}[
        every node/.style={font=\scriptsize\bfseries},
        box/.style={
          draw=stplred, line width=1pt,
          fill=stplcard, text=stplwhite,
          rounded corners=2pt,
          minimum width=2.2cm, minimum height=0.6cm,
          align=center,
        },
        arr/.style={-latex, draw=stplred, line width=1pt},
        lbl/.style={font=\tiny\itshape, text=stplmuted},
      ]
      \node[box] (wallet) {Portfel};
      \node[box, right=2.6cm of wallet] (net) {Walidatory};
      \node[box, right=2.6cm of net] (exp) {Explorer};
      \draw[arr] (wallet) -- node[lbl, above]{podpisana tx} (net);
      \draw[arr] (net) -- node[lbl, above]{finalized} (exp);
    \end{tikzpicture}
  \end{center}
  \vspace{0.15cm}
  \begin{exampleblock}{Co zobaczysz w Explorerze}
    Podpis (identyfikator), status i numer slotu, opłatę sieciową
    oraz listę kont, które zmieniły stan.
  \end{exampleblock}
\end{frame}

% TODO: 8b. Explorer --- pełnoekranowy zrzut transakcji na devnecie
\begin{frame}{Transakcja w Explorerze}
  \begin{center}
    \includegraphics[width=\linewidth,height=5cm,keepaspectratio]{images/explorer-tx-devnet-crop}
  \end{center}
  \begin{center}
    \textcolor{stplmuted}{\tiny explorer.solana.com · sieć: Devnet}
  \end{center}
\end{frame}

% ================================================================
\section{Podsumowanie}

% TODO: 9. Checklist --- gotowe do zadania blokującego
\begin{frame}{Checklist}
  \begin{enumerate}
    \item[{[\ ]}] Phantom zainstalowany w przeglądarce
    \item[{[\ ]}] Seed phrase zapisany offline
    \item[{[\ ]}] Sieć w portfelu ustawiona na \textbf{devnet}
    \item[{[\ ]}] Airdrop wykonany, saldo \textgreater{} 0 SOL
    \item[{[\ ]}] Pierwsza transakcja widoczna w Solana Explorer
  \end{enumerate}
  \vspace{0.5cm}
  \begin{center}
    \textcolor{stplmuted}{Masz wszystko do zadania blokującego.}
  \end{center}
\end{frame}

% ================================================================
\begin{frame}[plain]
  \begin{center}
    \vfill
    {\Huge\bfseries\color{stplwhite} gn}
    \vfill
  \end{center}
\end{frame}

\end{document}
